\section{Introduction}\label{sec:introduction}

For a binary word $x$ and a known deletion probability $q\in(0,1)$, let $\Dcal_q(x)$ denote the law obtained by deleting each bit independently with probability $q$ and concatenating the surviving bits in their original order.  The retention probability is $p=1-q$.  A sample from this law is a \emph{complete trace}: its length and all surviving bits are observed, without the positions they occupied in $x$.

For $s\in(0,1]$, let $T_{q,s}(n)$ be the least nonnegative integer $m$ for which an estimator, given $n,q$ and $m$ independent complete traces, recovers every $x\in\{0,1\}^n$ with probability at least $s$.  The estimator may be randomized and its computation is unrestricted.  If no finite $m$ suffices, set $T_{q,s}(n)=+\infty$.  Write $T_q(n)=T_{q,2/3}(n)$.  Our lower bounds concern this fixed sample budget.

\begin{theorem}[Quantitative sample lower bound]\label{thm:main}
Fix $s\in(0,1]$ and $0<c<1/(4\log2)$.  Along every sequence of instances $(n,q)$ with $n\to\infty$, $q\in(0,1)$, and $q^3\log n\to\infty$,
\begin{equation}\label{eq:main-bound}
 T_{q,s}(n)\ge n^{c\log(q^3\log n)}
\end{equation}
eventually.  All logarithms are natural.  In particular, for each fixed deletion probability $q\in(0,1)$ and fixed positive success probability $s$, the required sample number is $n^{\Omega(\log\log n)}$.
\end{theorem}

The varying-channel statement keeps the requested channel fixed throughout the construction for each instance; it is allowed to change only when the instance changes.  The result gives a negative answer to the polynomial-sample question for deterministic binary inputs at every fixed deletion probability.  A separate consequence concerns one-trace indistinguishability itself.  For laws on the countable set of finite binary words, write
\[
 \TV(P,Q)=\frac12\sum_w|P(w)-Q(w)|.
\]

\begin{theorem}[Superpolynomial one-trace indistinguishability]\label{thm:tv}
For every fixed deletion probability $q\in(0,1)$,
\begin{equation}\label{eq:min-tv}
 \lim_{n\to\infty}n^A
 \min_{\substack{x,y\in\{0,1\}^n\\x\ne y}}
 \TV\bigl(\Dcal_q(x),\Dcal_q(y)\bigr)=0
 \qquad\text{for every }A>0.
\end{equation}
At this same channel $T_{q,s}(n)/n^A\to\infty$ for every fixed $A>0$ and $s\in(0,1]$.
\end{theorem}

The channel is fixed independently of $A$.  The proof of \eqref{eq:min-tv} uses an explicit score and approximation bound.  Since the minimum is attained at every $n$, a single sequence of minimizing pairs realizes all the limits simultaneously.

\subsection{History and comparison with earlier results}

Reconstructing a sequence from repeated corrupted observations has a long history.  Levenshtein studied both combinatorial and memoryless channel formulations \citep{Levenshtein2001}.  Batu, Kannan, Khanna, and McGregor studied the independent-deletion trace problem in connection with multiple sequence alignment and outlined a linear sample obstruction at fixed deletion probability \citep[Section 4.2]{Batu2004}.  Later lower bounds for complete traces focused on the difficulty of locating a small discrepancy within a long background.  McGregor, Price, and Vorotnikova analyzed neighboring positions of a single one among zeros, using Hellinger distance and common mass \citep{MPV2014}.  Holden and Lyons used a defect in an alternating background to prove $\Omega(n^{5/4}/\sqrt{\log n})$ \citep{HoldenLyons2020,HoldenLyonsErratum2022}, and Chase obtained $\Omega(n^{3/2}/\log^7 n)$ by a direct Hellinger calculation for a related pair \citep{Chase2021Lower}.  The displayed Holden--Lyons and Chase rates fix the channel and a reconstruction success greater than $1/2$.  Our seed starts from the periodic defect family used by Chase and makes signed choices across many periods; it also supplies the overlap and score estimates needed for repeated amplification.

For fixed deletion probability, Holenstein, Mitzenmacher, Panigrahy, and Wieder obtained a worst-case upper bound of $\exp(\widetilde O(\sqrt n))$ \citep{HMPW2008}, where $\widetilde O$ suppresses polylogarithmic factors.  Independently, De, O'Donnell, and Servedio and Nazarov and Peres improved it to $\exp(O(n^{1/3}))$ \citep{DOS2017,NP2017}, and Chase obtained $\exp(O(n^{1/5}\log^5 n))$ \citep{Chase2021Upper}.  The 2026 preprint of Burudgunte, Valiant, and Wang gives
$\exp(p^{-7/3}(\log_2 n)^{C_0})$ samples for a universal constant $C_0$ and every positive retention probability $p=1-q$ \citep{BVW2026}.  At fixed $p$ this is a quasipolynomial sample upper bound; the cited preprint does not give a quasipolynomial-time reconstruction algorithm.  The logarithm of this sample bound is a fixed power of $\log n$, whereas the logarithm of our lower bound at a fixed channel has order $(\log n)(\log\log n)$, leaving a gap between the rates.  The companion paper \emph{Uniform quasipolynomial-time trace reconstruction} proves such a time bound for every fixed known rational retention probability \citep{OAI:Uniform-quasipolynomial-time-trace-reconstruction-October-5-2026}.

Polynomial reconstruction remains possible under other channel or input conditions.  For every fixed $\varepsilon>0$, Chen, De, Lee, Servedio, and Sinha give reconstruction in polynomial time when $q\le n^{-1/3-\varepsilon}$, using $O(n^{4/3})$ traces \citep{CDLSS2021LowDeletion}.  Those deletion probabilities lie outside the regime $q^3\log n\to\infty$ of Theorem~\ref{thm:main}.  Aamand, Liu, and Narayanan prove polynomial reconstruction when the internal zero gaps between successive ones have length at least a sufficiently large constant times $(\log n)^8$ and the deletion probability is at most a sufficiently small universal constant \citep{ALN2025}.

Some lower bounds restrict the available statistics.  De, O'Donnell, and Servedio measure the precision required of an oracle returning an $\ell^1$ approximation to the expected trace vector padded with zeros, while Nazarov and Peres study repeated observations of one fixed output coordinate \citep{DOS2017,NP2017}.  Chen, De, Lee, and Servedio study a tolerance oracle for functions of contiguous blocks of trace coordinates \citep{CDLS2024}.  Their lower bounds concern those specified statistics.  A different input model is studied by Rivkin, Valiant, and Valiant: the input is a fixed probability vector, with fresh Bernoulli bits drawn from it before each trace.  For deletion probability at least a constant times $n^{-1/2}$, their generalized trace laws can have total variation $\exp(-\Omega(\sqrt n))$ \citep{RivkinValiantValiant2025}.  Their arguments using translated segments and cancellation are related to the observation constructed in Section~\ref{sec:layout}.  Theorems~\ref{thm:main} and~\ref{thm:tv} concern the entire trace law of one deterministic binary word.

\subsection{Mechanism and proof structure}

Our construction hides differences in location at several scales.  The basic source of hidden randomness is exact common mass.  If trace laws for $k$ and $k+1$ copies of a word share mass $\omega$, place such a pair of runs on opposite sides of a central segment and choose their orientation by a fair bit.  Their joint trace laws share a product component of mass $\omega^2$.  On that component the orientation remains fair even after the side outputs are reported.  Many independent pairs therefore leave a binomial uncertainty in the position of the central segment.  Marginalizing the flags restores the product of the two side trace laws for each orientation.  Concatenating independent genuine traces of a partition of the fixed word then recovers its trace law exactly.  At later scales the preliminary observation uses earlier auxiliary outputs at its leaves and allows rare failed layouts; at every scale, conditioning to obtain likelihood domination adds a small change.  The resulting auxiliary laws therefore have an approximate common projection to genuine traces, with all these losses recorded in the invariant.

The two words use complementary signed patterns in their central segments.  The generating polynomial of a Prouhet--Thue--Morse pattern vanishes to high order at $1$ \citep[Proposition 2 and Theorem 6]{AlloucheShallit1999}, making its convolution with a broad binomial offset small.  The same dyadic sign recursion is used by \citet[Section 3.2, equation (3.10)]{HMPW2008} to produce iterated finite differences of a kernel for coordinate means.  In the conditional likelihood difference, each selected leaf set indexes a product of old likelihood scores, one per leaf, multiplied by a coefficient.  We expand this coefficient along the tree of nested segments as a sum of expected products of signed motif factors and support indicators, averaging over the hidden bits.  Within one product, a child is active at a vertex when it supplies one of these local factors, and the term is odd there when an odd number of the local factors are signed.  The complementary words give root types $+1$ and $-1$, so subtracting the two root evaluations cancels every even root term.  At a vertex where an odd term has only one active child, the hidden offset smooths its single local signed factor.  For an odd root term with no such singleton, selecting two active children recursively produces a forced binary subtree with $2^b$ selected leaves at depth $b$, and hence at least $2^b$ factors from the old score.  A shared displacement for each sibling group, together with independent local displacements for its members, spreads these descendant locations.  A joint bound that retains the shared displacement gives the extra decay beyond the factors from the old score.

The argument has three quantitative requirements.  The initial pair must work uniformly as $q$ decreases.  The number of amplification blocks grows with $n$, so constants depending on the channel or cancellation order must remain below the available scale margins.  Finally, every oracle conditioning and failed random layout must be charged to an error that remains negligible after taking many traces.

Section~\ref{sec:scales} states the oracle invariant and the global schedule with two phases.  Section~\ref{sec:tools} proves the common-mass, moment, and clipping tools, and Section~\ref{sec:initial} constructs the seed at the working channel.  Section~\ref{sec:block} chooses the local word and window sizes and proves the overlap estimates.  Section~\ref{sec:layout} defines the random observation, proves its projection to a genuine trace, and identifies the law of the hidden bits after the report.  Section~\ref{sec:coefficients} then proves the gain from the forced binary subtree and closes the induction.  Section~\ref{sec:transfer} restores genuine traces, pads to every exact length, and proves both main theorems.  The extension to arbitrary fixed positive success uses several independently chosen gadgets and posterior factorization.

\subsection{Distances and channel reduction}

For probability laws $P,Q$, define overlap by $\ov(P,Q)=1-\TV(P,Q)$, the largest mass of a common subprobability.  Our squared Hellinger distance is
\[
 H^2(P,Q)=\int(\sqrt{dP}-\sqrt{dQ})^2.
\]
Cauchy--Schwarz and Hellinger affinity give $\TV(P,Q)\le H(P,Q)\le\sqrt{2\TV(P,Q)}$.  Both distances decrease under a common stochastic kernel, and squared Hellinger distance of product laws is at most the sum for the factors.  If $Q\ll P$, then
\[
 H^2(P,Q)\le\int\left(\frac{dQ}{dP}-1\right)^2dP.
\]
These elementary inequalities also explain why a small likelihood score is an appropriate induction quantity.

For the rest of the proof put
\[
 \delta=\min(q,1/2),\qquad \Lambda=\delta^{-3},\qquad
 H_{\rm del}=\delta^3\log n.
\]
We use $\Tr(w)=\Dcal_\delta(w)$ throughout the construction.  The requested channel $q$ is obtained by deleting each surviving bit of a $\delta$-trace with additional probability $(q-\delta)/(1-\delta)$.  Thus all upper bounds on distinguishability proved at $\delta$ also hold at $q$.  Moreover
\[
 H_{\rm del}\to\infty,\qquad
 \log H_{\rm del}=\log(q^3\log n)+O(1).
\]
The working deletion rate $\delta$ remains fixed through every intermediate scale of an instance.  In all integer index ranges, endpoints may be real and only integer labels in the range are used.
